% ─────────────────────────────────────────────────────────────
% TSUE ANNOTATSIYA / ANNOTATION — TSUE BMI uslubiy ko'rsatma
% Section 5: max 300 words per language
% ─────────────────────────────────────────────────────────────

\newpage
\thispagestyle{plain}
\begin{center}
 {\large\bfseries ANNOTATSIYA}
\end{center}

\vspace{0.5em}
\noindent
Mazkur bitiruv malakaviy ishi O'zbekistonning 2017 yilgi iqtisodiy islohotlaridan so'ng Shimoliy-Sharqiy Osiyo (Xitoy, Yaponiya, Janubiy Koreya) ishlab chiqarish tarmoqlariga integratsiyalashuv darajasini tadqiq etadi. Tadqiqot O'zbekiston iqtisodiyotida tarkibiy asimmetriya mavjudligini isbotlaydi: mamlakatimiz oraliq mahsulotlar importi orqali keng ko'lamli orqaga integratsiya (backward integration) erishgan bo'lsa-da, ishlab chiqarilgan eksport tovarlar bo'yicha oldinga integratsiya (forward integration) sezilarli darajada qolib ketmoqda.

\textbf{Maqsad:} Shimoliy-Sharqiy Osiyo bilan savdo va to'g'ridan-to'g'ri xorijiy investitsiyalar (TXI) aloqalari ishlab chiqarish rivojlanishiga va iqtisodiy murakkablik ko'rsatkichiga qanday ta'sir ko'rsatishini empirik baholash.

\textbf{Metodologiya:} Tadqiqot Chow tarkibiy uzilish testini (2017 yil islohotlar nuqtasi sifatida), Balassa (1965) tomonidan ishlab chiqilgan oshkor qilingan qiyosiy ustunlik (RCA) indekslarini, ikki xil OLS regression modellarini va tarmoq tahlilini qo'llaydi. Ma'lumotlar Jahon banki (WDI), OEC BACI, UNCTAD va O'zbekiston Davlat statistika qo'mitasi bazalaridan olingan.

\textbf{Asosiy natijalar:} Chow testi 2017 yilni statistik jihatdan muhim tarkibiy burilish nuqtasi sifatida tasdiqladi ($F=6.937$, $p=0.013$), sanoat qo'shilgan qiymati o'sish sur'ati yiliga 0.4 dan 1.2 foiz punktgacha tezlashdi. Biroq TXI hajmi ($ \beta=0.002$, $p=0.15$) emas, balki oraliq mahsulotlar importi ulushi ($\beta=2.34$, $p<0.05$) sanoat modernizatsiyasining hal qiluvchi omili ekanligi aniqlandi. GVC boshqaruv tuzilmasi tahlili O'zbekistonning hali ham ``asir'' (captive) bosqichida ekanligini ko'rsatdi.

\textbf{Xulosa:} O'zbekiston ``asir integratsiya tuzoqidan'' chiqib, ``hamkorlik'' (relational) boshqaruv tuzilmasiga o'tishi uchun logistika to'siqlarini bartaraf etish, YaTO a'zoligini tezlashtirish va oraliq mahsulotlar savdosi chuqurligini oshirishga yo'naltirilgan sanoat siyosati zarur.

\vspace{0.8em}
\noindent\textbf{Kalit so'zlar:} \textit{global qiymat zanjiri, to'g'ridan-to'g'ri xorijiy investitsiyalar, O'zbekiston, Shimoliy-Sharqiy Osiyo, iqtisodiy murakkablik, sanoatlashish, tarkibiy uzilish, logistika, maxsus iqtisodiy zona.}

% ─────────────────────────────────────────────────────────────
\newpage
\thispagestyle{plain}
\begin{center}
 {\large\bfseries ANNOTATION}
\end{center}

\vspace{0.5em}
\noindent
This dissertation investigates Uzbekistan's integration into Northeast Asian (China, Japan, South Korea) production networks following the 2017 economic liberalisation reforms. It demonstrates a structural asymmetry at the core of Uzbekistan's industrialisation trajectory: the country has achieved significant backward GVC integration through the absorption of intermediate goods, but has failed to achieve commensurate forward integration into manufactured export markets.

\textbf{Research Objective:} To empirically assess the mechanisms through which trade and FDI relationships with Northeast Asia influence manufacturing value added and economic complexity.

\textbf{Methodology:} The study employs the Chow (1960) structural break test applied to manufacturing value added (2010--2023), Balassa (1965) Revealed Comparative Advantage indices, two OLS regression models with HC3 heteroskedasticity-consistent standard errors, directed trade network analysis, and a pooled country fixed-effects panel regression across five comparator economies.

\textbf{Key Findings:} The Chow test confirms 2017 as a statistically significant structural break ($F=6.937$, $p=0.013$), accelerating manufacturing growth from 0.4 to 1.2 percentage points annually. Critically, the FDI volume coefficient is statistically insignificant ($\beta=0.002$, $p=0.15$), while the intermediate goods import share is significant ($\beta=2.34$, $p<0.05$), identifying production network embedding rather than investment volume as the operative mechanism. GVC governance analysis confirms Uzbekistan remains in a captive rather than relational structure.

\textbf{Conclusion:} Escaping the captive GVC trap requires shifting from generic investment attraction toward policies that deepen intermediate goods trade relationships: supplier development programmes, backward linkage requirements in SEZ operating agreements, and WTO accession as an institutional lock-in mechanism.

\vspace{0.8em}
\noindent\textbf{Keywords:} \textit{global value chains, foreign direct investment, Uzbekistan, Northeast Asia, economic complexity, industrialisation, structural break, logistics, special economic zones.}
